\section{Тестирование пайплайна и модели}

\subsection{Подход к тестированию}

Тестирование построено на \textbf{синтетическом мини-датасете}, который генерируется фикстурой \texttt{tests/conftest.py} и полностью повторяет структуру реальных данных: недельные parquet-шарды взаимодействий, метаданные пользователей и видео, \texttt{npz}-файл эмбеддингов.

В синтетику намеренно закладываются:
\begin{itemize}
    \item \textbf{подсаженный сигнал}: каждому видео присваивается латентное качество $q$, от которого зависят и доля досмотра, и вероятности лайка/шеринга, а первая компонента эмбеддинга равна $q$.
    Благодаря этому агрегатные признаки предсказательны, и обученная на синтетике модель \emph{обязана} обыгрывать случайное ранжирование --- это проверяется тестом;
    \item \textbf{аномалии}: события с отрицательным и заведомо завышенным \texttt{timespent}, перекрученной долей досмотра --- для проверки очистки;
    \item \textbf{редкие сущности}: пользователь и видео с единственным событием --- для проверки фильтра минимальной активности.
\end{itemize}

Этапы пайплайна выполняются в тестовой сессии ровно один раз через цепочку session-scope фикстур:
\begin{center}
\texttt{pipe\_cfg} $\rightarrow$ \texttt{prepared} $\rightarrow$ \texttt{processed} $\rightarrow$ \texttt{featured} $\rightarrow$ \texttt{dataset} $\rightarrow$ \texttt{trained},
\end{center}
после чего отдельные тесты проверяют инварианты выходов каждого этапа.
Тесты не требуют ни GPU, ни сети, ни реальных данных.

\subsection{Уровни тестирования}

\textbf{Юнит-тесты чистых функций.}
Модуль метрик проверяется на ручных примерах (идеальное ранжирование дает $\text{NDCG}=1$; MRR и MAP сверены с расчетом на бумаге; группы без позитивов исключаются из усреднения) и статистической сверкой NDCG с эталонной реализацией scikit-learn на случайных данных.
Функции индексации данных проверяются на категоризацию файлов и временнýю сортировку шардов (случай \texttt{test/week\_26} лексикографически раньше \texttt{train/week\_00}).

\textbf{Интеграционные тесты этапов.}
Для каждого этапа проверяются содержательные инварианты, а не только отсутствие ошибок: формула таргета пересчитывается в тесте независимо и сравнивается поэлементно; агрегаты признаков пересчитываются вторым способом; проверяется, что события валидационного периода \emph{не участвуют} в агрегатах (утечка будущего) и что ни один компонент таргета не попал в список признаков.

\textbf{Smoke-тест обучения.}
На синтетике обучается настоящий \texttt{CatBoostRanker} (CPU, 60 итераций), после чего проверяются артефакты, диапазоны метрик и ключевое свойство: NDCG@10 и GAUC модели значимо превосходят случайное ранжирование (при этом GAUC самого случайного бейзлайна $\approx 0.5$ --- проверка корректности процедуры сравнения).

\subsection{Матрица покрытия}

Итоговое покрытие --- 40 тестов (таблица~\ref{tab:test-coverage}); полный прогон занимает $\sim$3 секунды.

\begin{table}[H]
\centering
\caption{Матрица покрытия тестами}
\label{tab:test-coverage}
\small
\begin{tabularx}{\textwidth}{@{}l l X@{}}
\toprule
\textbf{Файл} & \textbf{Тип} & \textbf{Проверяемые свойства} \\
\midrule
\texttt{test\_metrics.py} & юнит & NDCG/MAP/MRR/HitRate/Recall/Precision/GAUC: ручные примеры, сверка со sklearn, исключение групп без позитивов, диапазоны значений \\
\texttt{test\_prepare\_data.py} & юнит + интегр. & категоризация файлов, временнáя сортировка недель, состав manifest, ошибка при отсутствии каталога данных \\
\texttt{test\_preprocess.py} & интегр. & фильтры аномалий и активности, джойн метаданных, монотонность \texttt{event\_order}, поэлементное совпадение таргета с формулой, обнуление по дизлайку \\
\texttt{test\_features.py} & интегр. & совпадение агрегатов с независимым пересчетом, использование только train-окна, \texttt{emb\_cos} $\in [-1;1]$ и по строке на событие \\
\texttt{test\_build\_dataset.py} & интегр. & сплиты без пересечений и по границам времени, непрерывность групп, отсутствие компонент таргета в признаках, типы и пропуски \\
\texttt{test\_train.py} & smoke & артефакты обучения, структура и диапазоны метрик, превосходство над случайным ранжированием, важности всех признаков \\
\bottomrule
\end{tabularx}
\end{table}

\subsection{Запуск тестов}

\begin{lstlisting}[language=bash, caption={Запуск полного набора тестов}, label={lst:run-tests}]
python -m pytest tests/ -v
\end{lstlisting}

Результат прогона: \texttt{40 passed in 3.19s}.
Тесты также сыграли роль при переносе пайплайна на реальные данные: расхождения формата (эмбеддинги в \texttt{npz} вместо parquet, порядок недельных шардов) были оформлены как новые тестовые случаи после исправления.
